\subsection{Making Review Binding}
\label{sec:9.3.4}
This book has already said what an oversight board for sentient-AI research needs to look like; what makes its ruling something a researcher has to obey is a separate question, and human-subjects research already solved a version of it. The solution is institutional rather than criminal. An institutional review board's authority does not come from any power to fine or arrest a noncompliant researcher directly. It comes from a federalwide assurance: an institution promises, as a condition of eligibility for federal research funding, that all of its human-subjects research will go through review, and a federal office can suspend that assurance if the promise is broken — cutting off an entire institution's federal funding, not just the offending lab's \autocite{ohrp2018fwa}. The leverage is economic, not punitive, and it works because almost no research institution can function without the funding at stake.

The same structure transfers to sentient-AI research directly. A funding body, a journal's editorial standards, or a jurisdiction's research-licensing regime can condition eligibility on review-board clearance the same way federal grants already condition eligibility on IRB clearance — making noncompliance expensive to the institution rather than merely improper for the researcher. What the legal-status framework leaves open is which cases the board needs to see at all: a graded framework means low-capacity systems can be exempted from full review the way minimal-risk human studies already qualify for expedited review, freeing the board's scrutiny for the cases where a system's status is contested.

\runin{What the board applies}
The substantive standard is animal-research ethics' Three Rs. A binding board is what makes those three questions something a protocol has to answer before it runs. In practice, applying the Three Rs to a candidate sentient system means asking, in order: can a simulation or a non-sentient model answer the same question; if not, what is the minimum exposure that still produces a valid result; and given that minimum, what escalation triggers stop the study before harm compounds. A board with teeth can require a negative answer to the first question be documented, not merely asserted, before it will consider the second.

The Three Rs' own logic — defaulting toward the least invasive option available — argues for a specific evidentiary rule at every stage of this apparatus: a system's default experimental status stays the more protective one until the case for the less protective status is made, not the reverse. That default has a cost, and neither error it produces is free: it will sometimes extend protection to a system that turns out not to have needed it, and sometimes, before a threshold is met, withhold protection from one that did. The Three Rs answer only what a board does once a subject is inside its scope.
